\section{韩：小国的官僚化实验}

太行南口的关隘、通往洛阳的道路和散在山前的田邑，是韩国先要面对的疆土形状。韩氏从晋国裂解中取得的并非一块四面完整的领土，而是需要由宗族、家臣、城邑和道路接合起来的资源。前453年的\eventref{WQ-0453-ZHI}{战国·三家灭智}与前403年的\eventref{WQ-0403-THREEJIN}{战国·三家分晋}，让这种实际控制获得了新的名分；名分却不能替田邑收租，也不能替山口出守卒。约前448年前后，韩氏将所得田邑编入既有的宗族与家臣网络，经营太行南北的通道，\eventanchor{WQ-0448-HAN-JINYANG-ESTATES}{战国·韩氏经营分晋田邑}。这条国家线从一开始就有两面：道路能把粟帛、使者和兵员送到需要之处，敌对军队也能循同一条道路逼近。

前399年韩文侯即位，韩国仍须在郑、魏、赵之间经营这条狭长地带，\eventanchor{WQ-0399-HAN-WENHOU}{战国·韩文侯即位}。关于王年，《史记·韩世家》与《竹书纪年》的系统并非毫无异文；但把继承尽快接上，确是小国不让城邑、使聘和守备各自脱节的前提。约前四世纪初，韩国在郑韩旧城之间加固关隘与城防，\eventanchor{WQ-0392-HAN-DEFENSIVE-CITIES}{战国·韩国加固郑韩关隘}。这一判断可由《墨子》的守城议论和古城址调查照见，不能倒推为逐年工程簿册；它所显示的是道路一旦被夺，国内联络便可能断裂，守城所需的劳役和粮食也会长期压在城邑与乡里的身上。

账本所列前五世纪若干约年节点，把清点关隘、以田邑供役、集中粟帛、轮守山口和安置流动编户连成一组线索。它们不容许我们复原一份逐年颁发的韩氏簿册，更不能据此计算每座仓的存粮；但它们提示了韩国早期国家能力的实际材料。狭长土地难以长久供养一支能覆盖全部山口的常备军，守隘便要依赖田租、草料、修缮和轮番服役。这样的安排能够把分散资源变成短期防务，也会把劳役和转输的成本压回乡里，并使持久战争特别昂贵。

\begin{event}{前375年}{韩灭郑，迁都新郑}{可靠纪年}{郑地}\eventanchor{WQ-0375-HAN}{战国·韩灭郑}
《史记·韩世家》以极简的“因徙都郑”记下韩灭郑、迁都的结果。韩得到的不只是一片可在地图上涂色的土地，而是一座旧国都及其周边城邑、仓廪、市场、工匠和关津。把这些资源接入韩的征发与使聘网络，才会使吞并成为收入、兵员和调度能力；若接管失败，新得土地反会成为难以防守的裂缝。灭郑后接收新郑及周边城邑的吏治、仓廪，正是这一转换所不可少的一步，\eventanchor{WQ-0375-HAN-ZHENG-INTEGRATION}{战国·韩接收郑地吏治与仓廪}；接收的步骤没有连续记载，不能把它写成一份现存交接清册。新郑通向洛阳、陈地的道路与关津，因而改变了韩国交通的重心，\eventanchor{WQ-0375-HAN-ZHENG-ROUTE}{战国·韩接收郑地道路关津}，也使都城成为列国往来时很难绕开的地点。
\end{event}

新郑给韩带来了比旧日田邑更集中的都市资源。郑韩故城的考古与传世叙事可以相互照亮这种城市性：接收郑都市场、工匠与交通税源，强化了新郑的都城功能，\eventanchor{WQ-0375-HAN-ZHENG-MARKET}{战国·韩接收郑都市场与工匠}，却不能替我们开出一张战国税额表；城墙、仓廪和道路可以支撑守备，也需要持续的劳力与粮食。韩所取得的是一处可以组织国家行动的中心，而非一座足以隔绝风险的堡垒。西面秦东出的压力、北面魏赵的竞争、南面楚的影响，都会在这片交通网中留下痕迹。新郑的繁荣和暴露，正是同一地理条件的两种后果。

前377年前后韩哀侯即位，\eventanchor{WQ-0377-HAN-AIHOU}{战国·韩哀侯即位}；灭郑后不久，他又于前374年去世，韩懿侯继位，\eventanchor{WQ-0374-HAN-YIHOU}{战国·韩懿侯继位}\eventanchor{WQ-0374-HAN-AIHOU-DEATH}{战国·韩哀侯去世与郑地重整}。世系年次在《史记·韩世家》与《竹书纪年》中有异文，宫廷又如何逐一安置新旧贵族、郑地官属和各地命令，材料不足以重演。前363年韩懿侯卒、韩昭侯继位，\eventanchor{WQ-0363-HAN-ZHAOHOU}{战国·韩昭侯即位}，才为后来的用人和整顿提供了较稳定的君主支点。对一个刚把旧国都纳入本国的政权而言，继承的连续性本身就是行政与守备的条件；失去它，仓廪、关津和城防便可能重新为不同的政治网络所用。

\begin{event}{约前354年至前337年}{申不害相韩}{王年较可靠，制度细节有争议}{韩}\eventanchor{WQ-0354-SHENBUHAI}{战国·申不害相韩}
《史记·韩世家》记申不害相韩昭侯十五年；《韩非子·定法》则把昭侯用申不害后的韩写成“国内以治，诸侯不来侵伐”。两种材料都保存了韩国在强邻间整顿内部的记忆，但后一句出自战国末年以法、术、势论政的著作，不能当作相府留下的逐日政绩。较稳妥的结论是，韩廷试图让官吏所担的名目同实际完成的事务可以相校，让任官、赏罚和命令不至全被宗族、私交或地方关系截断。
\end{event}

后世以“术”概括申不害，常将循名责实、考核官吏和君主控臣联结在一起。放回新郑的处境，它首先不是一套悬在田地、仓廪和关隘之外的玄理。君主若不能较可靠地知道谁应征发、谁管守备、谁负责转输，有限的收入和人力就会在层层关系中耗散；相反，官名、功绩与责任若可被核验，才较可能把城邑的户口、田租、道路和军粮接到同一条命令链上。约前350年至前340年的名实考核努力，\eventanchor{WQ-0345-HAN-SHENBUHAI-ADMIN}{战国·申不害推行名实考核}，正指向这种缩短命令与结果之间距离的尝试。

但这不等于申不害已颁行一部可以逐条复原的行政法典。其人著作早佚，今见《申子》多赖辑佚；韩非对“术”的解释也服务于他自己的论证。约前305年前后，申不害去世后考核官吏与君主控臣的政治语言仍在韩国流传，\eventanchor{WQ-0305-HAN-SHENBUHAI-LEGACY}{战国·申不害术治的后续遗产}；后世法家文本所概括的刑名、簿书和属官考核，也提示宫廷仍以此维持控制，\eventanchor{WQ-0305-HAN-SHENBUHAI-REGISTERS}{战国·韩国延续簿书与官吏考核}。这些并非韩国档案，不能倒推每项清核户口、仓储田租、道路传递、田宅或文书责任的办法皆为申不害亲定，更不能断言它们在全国以同一格式执行。行政技术能够使资源较少漏损，也会使邑吏和编户承受更细密的登记、差役与问责；它改善的是调度的可能，而不是凭空增加土地、人口或安全纵深。

\begin{sourcenote}[申不害的“术”：理论遗产不是相府档案]
\sourcebook{史记·韩世家}提供相韩的基本政治线索；\sourcebook{韩非子·定法}保存后人如何把申不害置于法、术、势的争论中；辑佚的\sourcebook{申子}只能补出零散概念。三者足以支持“名实相校、考察官吏”是韩国政治传统的重要线索，却不足以证明某一官署的编制、每一次考课或一段劝说原话。有关城市、道路和仓廪的考古与历史地理研究，可以帮助追问制度所依托的物质条件，不能替代韩国官府的直接文书。
\end{sourcenote}

申不害式的整顿也没有让韩摆脱同为三晋的魏。前343年至前342年，魏攻韩，韩遣使向齐求援，\eventanchor{WQ-0343-WEI-HAN}{战国·魏攻韩与齐援}。这不是一则小国以巧言转危为安的故事：韩国单独抵抗的能力有限，向齐求援会把韩魏冲突牵入齐魏的力量竞争；齐的介入最终转向\eventref{WQ-0342-MALING}{战国·马陵之战}的战局。传世材料保存了求援和齐魏对抗的基本轮廓，却不足以逐字还原韩使的陈词，或将韩的选择解释为预先算定的谋略。对于新郑而言，这次危机表明，行政上的可调度并不等于能独自决定战争的地点和援军到达的时日。

申不害退出政治舞台以后，君位更替仍不断把韩廷置于战线旁。前312年韩宣惠王卒、韩襄王继位，\eventanchor{WQ-0312-HAN-XUANHUI}{战国·韩襄王继位}；这件事发生在秦、楚冲突加剧之时，具体的宫廷会议无从复原。前308年，秦军久攻宜阳，须从关中、河西转运粮食、器械并轮换士卒，韩军则依山城迟滞战局，\eventanchor{WQ-0308-YIYANG-LONG-SIEGE}{战国·秦久攻宜阳}。《史记·樗里子甘茂列传》所见围城时长及细节还须校勘版本，不过宜阳控制洛阳西南山口、攻方必须同时维持转输与君臣信任的难题，并不因此消失。

这也说明“内治”与“外患”不能被截然分开。户口、仓储、道路和文书之所以值得清核，正因为一旦魏军压境，守城、野战、使聘和粮道会同时索取这些资源；而这些资源被集中到新郑后，外围城邑便未必能得到同样的保护。韩昭侯以后，申不害留下的考核语言仍可被后继士人和官府使用，可它无法替代一支拥有更大田赋、更稳定后方或更可信外援的军队。

\begin{event}{前300年至前293年}{西部城邑失守与伊阙受创}{可靠纪年；战果细节有争议}{秦、韩、魏}
前300年秦取新城，\eventanchor{WQ-0300-QIN-XINCHENG}{战国·秦取新城}，前294年向寿又取武始，\eventanchor{WQ-0294-XIANGSHOU-WUSHI}{战国·向寿取武始}，韩国西部城邑被逐点切割；两城的地望仍有考证问题，不能据地名画出精确的推进线。前293年，韩魏集结军队于伊阙，试图守住伊、洛之间的狭谷和两国交通线，\eventanchor{WQ-0293-HAN-WEI-YIQUE-COALITION}{战国·韩魏会师伊阙}，却在\eventref{WQ-0293-YIQUE}{战国·伊阙之战}中遭到决定性挫败。秦军可以把关中方向的兵员和粮食推向东面；韩、魏却须在越来越窄的城链、渡口和山口之间分配守军。战后传称斩首二十四万，数字在《史记·白起王翦列传》与后来的战功叙事中格外醒目；它不能当作人口统计或可核对的伤亡表。无需倚赖这一巨数，韩魏的野战能力和西向屏障受重创，已足以解释韩国为何此后更常把割地、守城和求援并用。
\end{event}

韩襄王继位之际所面对的夹击并非只有一条西线：魏仍是近处竞争者，秦的东进日益直接，楚在南方的行动也会改变中原联盟的价格。对韩廷而言，结交强邻可以换取一时的缓冲，却很难把自己的城邑从别国的行军图上抹去。伊阙之后，韩国须重新衡量洛阳西南城链、保存新郑核心，并把割地与求援并用，\eventanchor{WQ-0293-HAN-YIQUE-REASSESSMENT}{战国·伊阙后韩国重整西部防线}。前290年向秦割让武遂等地求和，\eventanchor{WQ-0290-HAN-WUSUI-CESSION}{战国·韩割武遂等地}；地名、辖属及具体边界虽有异说，防线向太行南口和新郑西北后退这一趋势却很清楚。

前274年前后，新郑优先整修仓廪和城门，保护王畿粮道，\eventanchor{WQ-0274-HAN-NEWZHENG-GRAIN}{战国·韩整修新郑仓廪与城门}；又修补西北山口城塞，试图以分层防线防止秦军切断王畿道路，\eventanchor{WQ-0274-HAN-NEWZHENG-PASSES}{战国·韩修补新郑西北山口}。这些工程的年次、仓廪位置和施工量都没有完整文书，不能被涂成一项万全的“国防计划”。不过，城防与仓储的结合揭出韩国的现实尺度：把粮食、征发和守军收束在都城周边，能够增加危急时的承受力；代价是周边城邑更容易在秦军逐步推进时被单独切离。行政的精密在这里并非军力的替身，而是以有限土地和税源维持有限防线的办法。

\begin{event}{前273年}{华阳援韩与联军失利}{可靠纪年；联军指挥细节有争议}{韩、赵、魏、秦}\eventanchor{WQ-0273-HAN-HUAYANG-RESPONSE}{战国·华阳援韩}
秦攻韩时，韩国向赵、魏求援，并调集新郑周边守军；赵魏援军在华阳为秦所败。三国都有阻止秦继续东进的理由，却未因此拥有一套可由韩国指挥的共同军令、补给和进退次序。现存《史记》《战国策》及《资治通鉴》的材料能确认援韩、交战与失利的主线，不能把联军的指挥结构、每一方的先后到达或劝说之辞补写成完整战报。
\end{event}

华阳使韩的困境格外清楚：外交不是同军政相分离的一门口才，而是把别国的兵、粮和行动时间接到本国危机上的艰难工作。新郑可以动员本地守军，却无法保证赵魏在同一地点、同一时刻承担同样的风险；赵魏也要计算自己的边境和收获。韩釐王在这一年去世，韩桓惠王继位，\eventanchor{WQ-0273-HAN-XI-DEATH}{战国·韩桓惠王继位}，所面对的已是更依赖割地和短暂外交缓冲的国家。把这种收缩归结为某一位君主或使者的胆怯，既超出材料，也会忽略华阳所显示的联盟协调成本。

\begin{event}{前264年至前262年}{上党交通断裂}{可靠纪年；地方决策动机有争议}{上党、韩、赵、秦}\eventanchor{WQ-0264-HAN-SHANGDANG-CRISIS}{战国·上党危机}
白起攻陉城等要点后，韩国以太行山口城塞分散阻滞秦军，\eventanchor{WQ-0264-HAN-MOUNTAIN-PASSES}{战国·韩以太行山口阻秦}，可以延缓进逼，却难以恢复同上党的安全联络。上党守吏随之加强壶关、长子等山口防务，\eventanchor{WQ-0264-HAN-SHANGDANG-PASSES}{战国·上党加强山口防务}；各城失守先后、守吏姓名和陉城地望皆有争议，不能把这些布置排成完整战报。到前262年，秦五大夫贲攻韩取十城，\eventanchor{WQ-0262-QIN-BEN-TEN-HAN-CITIES}{战国·秦五大夫贲取韩十城}，\eventref{WQ-0262-SHANGDANG}{战国·上党归赵}随之发生。上党守吏在与新郑交通断绝的条件下拒绝降秦、转向赵求援，使一场韩秦之间的边地危机升级为秦赵争夺。后出的叙事给这项选择安放了鲜明的言辞和人物动机；可确知的，是道路断裂、守备压力和赵的介入共同改变了战场的尺度。
\end{event}

上党并非韩国可以随手舍弃的一角，也不是一项行政“术”本该自动解决的失误。它把太行山口、城邑收入和军事指挥线系在一起：一处要点失守，地方守军不仅失去援兵，也失去同都城核对命令、转送粮秣和请求支援的条件。山口防线能够争取时间，不能制造纵深；城邑越被逐点夺取，原先可由新郑统一调度的租赋和兵役越少。上党转归赵把韩国的局部困境推入秦赵之间的大战，随后\eventref{WQ-0260-CHANGPING}{战国·长平之战}又重塑整个中原的均势；这条链条有明确的触发因素，却不该被写成从伊阙到统一的一条自动直线。

长平之后，韩仍有朝廷、城市和政治选择，并未立刻消失。它保住新郑的努力，也有自己的社会代价：每一次清核户口、重排仓储田租、整理道路渡口和复核差役，都是要把减少了的资源重新转成守备与行政。账本在前290年、前234年和前230年所列的这些动员性节点，宜理解为持续的治理难题，而非已经发现的逐条诏令。它们显示，韩国末期并非全然停止治理；恰恰因为核心区仍须运转，编户、邑吏与承担运输和差役的人便要承受更密集的压力。

\begin{event}{约前260年至前233年}{韩非的著述与入秦}{年代与文本层次有争议}{韩、秦}\eventanchor{WQ-0234-HANFEI-QIN-RECEPTION}{战国·韩非入秦}
韩非在韩国危机中撰述《孤愤》《五蠹》等法术思想文本，\eventanchor{WQ-0260-HANFEI-WRITING}{战国·韩非撰述法术文本}，后为秦廷所读并被召入秦。篇章的写作年次与《韩非子》的成书层次仍有争议，不能把书中诊断当作某一年新郑政策的记录。韩国试图利用本国公子的思想声望争取缓冲，所面对的却是一个军事上已经难以单独相持的对手。关于韩非入秦后遭李斯等构陷而死的传记叙事，见于《史记·老子韩非列传》《李斯列传》，人物关系和言辞经过强烈的后出塑形，不能当作秦廷密谈的逐字记录。可以确认的结果是，韩非未能改变秦的吞并路线；他的法术论述却在秦汉以后的政治思想中被不断重读。
\end{event}

韩非使韩国的困境留下了最有影响力的理论语言，却不能把理论倒当作韩国的实际政令。《韩非子》把君主如何防止臣下专权、怎样审察名实、如何用法与势组织国家推到极端清楚的程度；它既吸收申不害传统，也以战国末年士人的问题意识重组了它。文本批评所见的成书层次、传记的戏剧化以及韩非具体活动的稀少记录，都要求我们不把《孤愤》中的诊断直接等同于新郑的政策，更不能把书中的说辞当作韩廷、秦廷的实录。韩非的重要性正在于此：他让一个受挤压国家对于官僚、君权和外交困局的反思获得了远超过韩国疆域的生命。

前254年，韩桓惠王入秦朝见，\eventanchor{WQ-0254-HAN-HUANHUI-COURT-QIN}{战国·韩桓惠王朝秦}。这至少显示韩以君主亲至承认秦的军事优势、换取暂缓直接吞并的压力；朝见礼仪、割地和口头承诺没有完整条约，不能补成确定的屈服仪式。前239年桓惠王卒、韩王安继位，\eventanchor{WQ-0239-HAN-AN}{战国·韩王安继位}。新君在秦与赵魏之间维持使聘、以从属外交延缓新郑危机，\eventanchor{WQ-0239-HAN-AN-FOREIGN-POLICY}{战国·韩王安维持从属外交}，但早期具体使节活动记录稀少；能够确定的只是，韩国已经难以用独立作战改变秦先取近邻的战略排序。

\begin{event}{前231年至前230年}{献南阳与韩国覆亡}{可靠纪年；接管细节有争议}{南阳、新郑}\eventanchor{WQ-0231-HAN-NANYANG-HANDOVER}{战国·韩献南阳}
前231年，韩国献南阳，秦以腾治其地，并接手城邑、道路和仓储，\eventanchor{WQ-0231-HAN-NANYANG-ROADS}{战国·秦接管南阳道路与仓廪}。对韩而言，割地或许能换得片刻存续；对秦而言，取得的是直通新郑的行政前沿。内史腾的官号、辖区、南阳范围及接管的细部仍有讨论，不能据此想象一份保存完好的交接清册；但前沿一旦由秦官经营，韩国核心区便更直接暴露在秦的军政网络之前。前230年，\eventref{QIN-0230-HAN}{秦·灭韩}，韩王安被俘，\eventanchor{WQ-0230-HAN-AN-CAPTIVITY}{战国·韩王安被俘}，韩国作为独立政权终结；旧臣在秦郡县任职或离散的具体路径，史料只留下零星痕迹。
\end{event}

这不是一座城在某一日忽然失守便能说明的结局。《史记·秦始皇本纪》与《六国年表》给出灭韩的可靠年代，却没有留下可逐日复原的攻城、降服和行政接收记录。长期的结构性原因，是狭长国土必须同时供给城防、道路、朝廷和多方向外交，秦的持续东进则不断缩短其可供调度的空间；献南阳和前230年的进军是使这种脆弱性成为终局的直接节点。行政考核曾帮助韩把新郑、田地和人力组织得更紧，却不能在失去上党、西部城邑和南阳之后重新长出收入、同盟或纵深。

韩国的历史也不应在此被写成秦统一的序言。前226年，韩王安被迁处，传称后为秦所杀，\eventanchor{WQ-0226-HAN-AN-REMOVAL}{战国·韩王安被迁处}，显示旧王族仍是统治者须处置的政治象征；其结局年份和被杀说皆属后出纪事，不能坐实。前218年前后，张良等旧韩贵族联结宾客、力士和地方交通组织反秦，\eventanchor{QIN-0218-HAN-REMNANT-NETWORK}{秦·旧韩贵族联结反秦网络}；博浪沙事件后，秦又加强对旧韩宗族、宾客与道路往来的监视，\eventanchor{QIN-0218-HAN-ARISTOCRAT-SURVEILLANCE}{秦·监视旧韩宗族与宾客}。这些事主要见于《史记·留侯世家》《刺客列传》与《秦始皇本纪》，有关网络的规模、密谋的细部、器械和人物言辞都不能随意补足。它们至少说明，新郑周边的道路、旧贵族的声望和私人联系并未随着国家名号消失；曾经维系小国存续的社会资源，在帝国接管之后仍可能成为不服从的条件。

韩国的编年中还有一些看似只是王位或城邑的短条，却使上述过程不至悬在空中。韩文侯即位后经营郑、魏、赵之间的狭长土地，\eventref{WQ-0399-HAN-WENHOU}{战国·韩文侯即位}；韩哀侯、懿侯、昭侯的继承，\eventref{WQ-0377-HAN-AIHOU}{战国·韩哀侯即位}\eventref{WQ-0374-HAN-YIHOU}{战国·韩懿侯即位}\eventref{WQ-0363-HAN-ZHAOHOU}{战国·韩昭侯即位}，与前375年接收郑地官属、市场和道路的任务交叠在一起。没有现存韩廷簿册能让我们量出新郑市场的税额或每处关隘的守卒；郑韩故城和《韩世家》足以说明，都城接收从来不是一日完成的改名。每一次继承都要让新旧官属继续承认命令，正是小国行政的脆弱起点。

申不害之后，宜阳久攻、伊阙集结与武遂割让，把这种脆弱性推到山口和粮道上，\eventref{WQ-0308-YIYANG-LONG-SIEGE}{战国·宜阳久攻}\eventref{WQ-0293-HAN-WEI-YIQUE-COALITION}{战国·韩魏守伊阙}\eventref{WQ-0290-HAN-WUSUI-CESSION}{战国·武遂割让}。它们不能证明申不害留下了一套全境有效的簿书机器；反而说明即使有考核与文书，韩国仍须在战时修补新郑仓廪、太行山口与上党通道，\eventref{WQ-0274-HAN-NEWZHENG-PASSES}{战国·新郑修山口}\eventref{WQ-0264-HAN-SHANGDANG-PASSES}{战国·上党守隘}。前254年韩王入秦朝见、前239年韩王安继位，\eventref{WQ-0254-HAN-HUANHUI-COURT-QIN}{战国·韩王朝秦}\eventref{WQ-0239-HAN-AN-FOREIGN-POLICY}{战国·韩王安即位}，并不是“屈服”一词可以替代的外交：它是以君主、使聘和剩余城邑换取不确定喘息的连续试探。

\begin{turningpoint}[制度得失：能收束资源，不能取消位置]
韩国的实验解决过真实问题：把新郑的城市资源、田邑收入、官吏责任、道路转输和守备需要较紧地接在一起，使小国不必完全依赖旧宗族的私属网络。承担成本的是被清点、征发、轮守和转输的编户与邑吏；当战线缩短时，负担还会更集中到核心区。申不害及其后的行政语言并非无效，正因为它能提高有限资源的可用性；但国家能力从来受资源规模、地理纵深、盟军协调和敌方持续投入的边界限制。韩的覆亡因此不是“术”徒劳的判决，而是一个官僚化小国在列强间把可动员之物用到极限后，仍无法独自决定其安全环境的历史。
\end{turningpoint}
